\subsection{Why placement is difficult}\label{sec:background}

A netlist specifies circuit components and the pins that must be connected.
Physical design assigns locations to those components and routes wires between
their pins on silicon. Hard macros, such
as SRAM blocks, have fixed dimensions and pin locations relative to the block;
the surrounding logic is implemented with much smaller standard cells. Macro
placement chooses where the large blocks sit, leaving room for the cells and
their connections. Later routing stages choose metal tracks and vias. Placement
therefore shapes a routing problem whose final paths are still unknown
\citep{openroadflow}.

The decisions are coupled through both nets and geometry. Moving a macro toward
one neighbor can lengthen connections to others and displace demand into a
different part of the canvas. IBM01 alone has 246 hard macros,
894 soft standard-cell clusters, and 5,993 loader nets with at least one resolved
benchmark node. A useful
location depends on the rest of that arrangement, not just the macro's own size
or strongest connection.

Wirelength and available routing space also compete. Bringing connected pins
together shortens their distance, but crowding blocks can leave narrow channels
for other nets. Spreading blocks can open those channels while lengthening
connections between the blocks, as Figure~\ref{fig:placement-tradeoff}
illustrates. Local density measures how much area is occupied; congestion
measures wiring demand relative to routing resources. A region can have spare
placement area yet attract too many connections. Routability-driven placers
therefore account for more than wirelength \citep{cheng2019replace,openroadmacro}.

Legal motion creates a second difficulty. Two hard rectangles must be separated
horizontally or vertically. A nearby improvement in their current ordering may
be easy to find, while exchanging their roles can require several blocks to move
together or an intermediate increase in cost. A one-coordinate search tests
only a slice of this space. The continuous coordinates and discrete relative
arrangements must be handled together; Section~\ref{sec:geometry-analysis}
describes the resulting feasible set.

Finally, each candidate has to be evaluated. Pin extrema, grid intersections,
and selection of the most loaded grid bins make the proxy piecewise rather than
everywhere smooth. Full routing is a later validation step; the challenge uses
the proxy to guide search within a one-hour limit per design \citep{challenge2026}.
Even full proxy recomputation took seconds on the measured IBM01 and IBM10
move sequences, making the cost of repeated trials consequential.
Thinkorplace uses smooth models for coordinated proposals, legalization to
restore hard geometry, and incremental scoring to make many local trials
affordable. These choices address different parts of the problem: finding a
promising arrangement, reaching a feasible one, and checking whether a move
improves the stated objective.

\begin{figure}[htbp]
\centering
\begin{tikzpicture}[x=0.73cm,y=0.53cm,font=\small,
  block/.style={draw=blue!65,fill=blue!8,line width=0.6pt},
  route/.style={draw=orange,line width=0.7pt},
  connection/.style={draw=blue,line width=1.3pt}]
  \begin{scope}
    \node[font=\small\bfseries] at (3,6.55) {Compact arrangement};
    \draw[rulegray] (0,0) rectangle (6,6);
    \path[fill=orange!10] (0,2.65) rectangle (6,3.35);
    \draw[block] (0.9,3.35) rectangle (5.1,4.85);
    \draw[block] (0.9,1.15) rectangle (5.1,2.65);
    \node at (3,4.1) {Macro A};
    \node at (3,1.9) {Macro B};
    \foreach \y in {2.8,3.0,3.2} {\draw[route] (0.15,\y) -- (5.85,\y);}
    \draw[connection] (3,2.65) -- (3,3.35);
    \fill[blue] (3,2.65) circle (1.5pt);
    \fill[blue] (3,3.35) circle (1.5pt);
    \node[align=center,font=\footnotesize] at (3,-0.72)
      {Shorter A--B connection;\\narrower channel};
  \end{scope}
  \begin{scope}[xshift=5.6cm]
    \node[font=\small\bfseries] at (3,6.55) {More space between blocks};
    \draw[rulegray] (0,0) rectangle (6,6);
    \path[fill=orange!10] (0,2.0) rectangle (6,4.0);
    \draw[block] (0.9,4.0) rectangle (5.1,5.5);
    \draw[block] (0.9,0.5) rectangle (5.1,2.0);
    \node at (3,4.75) {Macro A};
    \node at (3,1.25) {Macro B};
    \foreach \y in {2.4,3.0,3.6} {\draw[route] (0.15,\y) -- (5.85,\y);}
    \draw[connection] (3,2.0) -- (3,4.0);
    \fill[blue] (3,2.0) circle (1.5pt);
    \fill[blue] (3,4.0) circle (1.5pt);
    \node[align=center,font=\footnotesize] at (3,-0.72)
      {Longer A--B connection;\\wider channel};
  \end{scope}
  \draw[connection] (0,-1.85) -- (0.7,-1.85);
  \node[anchor=west,font=\footnotesize] at (0.85,-1.85) {A--B connection};
  \draw[route] (5,-1.85) -- (5.7,-1.85);
  \node[anchor=west,font=\footnotesize] at (5.85,-1.85) {Other nets crossing the channel};
\end{tikzpicture}
\caption{The same two blocks in two legal placements. Increasing their
separation leaves more room for horizontal routing but lengthens the vertical
A--B connection. The illustration assumes a routing layer obstructed by the
macro interiors. Colored crossings represent separate routing layers without
electrical junctions. This is a schematic tradeoff, not a measured benchmark.}
\label{fig:placement-tradeoff}
\end{figure}
